\section{From Forecast to Trading Decision}

\subsection{Why forecasts must feed decisions}

A price forecast is not the end product. It is an input to a decision.
Understanding that decision is what separates a forecasting engineer
from a forecasting scientist.

On a trading desk, the decision is: how much power to buy or sell
in the day-ahead auction, at what price.
A forecast helps if it changes what the desk does.
If the forecast is the same as the TSO's, the desk learns nothing new.

\subsection{The bid optimisation problem}

A balance-responsible party (BRP) must submit bids for each of the 24
hours of tomorrow. For each hour $h$, they choose:

\begin{itemize}
  \item \textbf{Quantity} $q_h$ (MWh): how much to buy/sell.
  \item \textbf{Price} $p_h$ (EUR/MWh): the limit price (below which they do not execute).
\end{itemize}

If they own generation, they want to sell when price is high.
If they own consumption, they want to buy when price is low.
If they get the timing wrong, they pay the balancing market.

A better price forecast lets them time bids more accurately.

\subsection{Expected value of forecast improvement}

Suppose the desk has 100 MW of flexible consumption (a factory).
Daily energy cost = $\sum_{h=1}^{24} q_h \cdot p_h^{\text{actual}}$.

Naive strategy: buy flat (4 MWh/h = 100 MW $\times$ 1h). No timing.
Smart strategy: shift consumption to the cheapest 8 hours.

Over 2 years of PL data, the spread between cheapest 8 and most expensive 8
hours averages 47 EUR/MWh. At 100 MW that is 4,700 EUR/day, 1.7M EUR/year.
A 10\% improvement in price forecast accuracy is not nothing.

\subsection{Where probabilistic output matters most}

P50 (median) finds the best hour to buy.
P10 and P90 answer: \emph{how confident are we?}

A trader will not shift 100 MW of load if the P90 suggests prices
might still be high. The decision threshold depends on the band width.

Wide band: stick closer to the baseline (less timing risk).
Narrow band: push harder on the optimal hour.
The quantile forecast is the input to a \emph{risk-adjusted} decision.

\subsection{Value at Risk and forecast uncertainty}

A portfolio optimiser uses price scenarios, not point forecasts.
Typical workflow:

\begin{enumerate}
  \item Sample $N=1000$ price paths from the P10/P90 band.
  \item For each path, compute portfolio value.
  \item Report Value at Risk: the 5th percentile of portfolio value across paths.
  \item Accept bids only when VaR is above the position limit.
\end{enumerate}

This is why the CQR upgrade matters: if the band is too narrow,
VaR is underestimated, positions are too large, losses can exceed limits.

\subsection{The merit-order link}

Price is set by the most expensive running unit (section 15).
A good price forecast is therefore also an implicit capacity forecast:

\begin{itemize}
  \item Low solar + high demand → gas marginal → 90--120 EUR/MWh.
  \item High wind + low demand → wind/coal marginal → 30--60 EUR/MWh.
  \item Emergency: must-run capacity fails → reserve scarcity pricing → 300+ EUR/MWh.
\end{itemize}

Our LGBM SHAP says solar forecast is driver \#1 (mean $|\text{SHAP}|$ = 18.7 EUR/MWh).
That is the merit-order link made quantitative.

\subsection{Imbalance settlement}

After gate closure, reality deviates from the declared position.
The deviation is settled at the balancing price (section 10).

Since June 2024 PSE reform: balancing uses scarcity pricing.
When the system is short, the balancing price can exceed 500 EUR/MWh.
When long (oversupply), balancing price drops below day-ahead.

A BRP with a good load forecast minimises deviations.
Each 1 MWh deviation at 100 EUR imbalance price = 100 EUR cost.
At 100 MW peak, even 1\% MAPE improvement saves meaningful money.

\subsection{Capacity market connection}

Rynek mocy (section 11) contracts pay capacity providers
to guarantee availability, regardless of how much they generate.
A capacity auction 5 years ahead requires forecasts of:

\begin{itemize}
  \item Future peak demand (load model).
  \item Future renewable generation profiles (reducing residual demand).
  \item Whether a given unit will be called as the marginal unit.
\end{itemize}

The day-ahead price model provides the market context:
if solar growth continues to suppress midday prices,
midday peakers lose revenue → they exit → future peak scarcity risk rises.

This is the long-run equilibrium that capacity markets try to prevent.

\subsection{Interview lines}

\begin{itemize}
  \item ``The value of a price forecast is in the decisions it improves.
        On a 100 MW flexible load, a 1 EUR/MWh MAE improvement is worth
        roughly 365 EUR/day in expected savings. That is the conversation
        I would have with a trading desk.''
  \item ``Probabilistic output is not a nice-to-have. It is the input to
        risk-adjusted decision-making. P50 alone cannot size positions.''
  \item ``Our SHAP analysis shows solar forecast is price driver \#1.
        That means the quality of the RES forecast limits price forecast
        quality — a direct link from meteorology to trading P\&L.''
  \item ``Imbalance settlement since June 2024 uses scarcity pricing.
        The payoff to load forecast accuracy has increased.''
\end{itemize}
